% ============================================================
%  RESULTS — Master Thesis
%  "Offensive Corner Tactics vs. Defensive Structures:
%   An Automated Classification and Effectiveness Analysis"
% ============================================================

\chapter{Results}
\label{ch:results}

This chapter reports the empirical findings produced by the analytical
pipeline specified in Chapter~\ref{sec:methods}. Findings are presented in
the order of the sub-questions: SQ1 (defensive setup classification),
SQ2 (offensive pattern classification), SQ3 (effectiveness of offensive
strategies against defensive structures), and SQ4 (individual player
qualities). The applied integration of these results into the Data \&
Video Analyser (SQ5) is illustrated through a single worked example at
the end of the chapter; the design and data contract of the integration
itself was specified in Section~\ref{sec:sq5_dv_analyser}.

Throughout, results are presented with explicit attention to sample size
and to the data-coverage limitations identified in
Chapter~\ref{sec:methods}. A finding based on a small cell or on the
86-corner manually labelled subset is flagged as such where it appears,
and is interpreted only as exploratory unless supported by the full
1{,}401-corner dataset.

\section{SQ1: Defensive Setup Classification}
\label{sec:results_sq1}

\subsection{Individual defender role classification}

The ensemble classifier (Random Forest, Extra Trees, and Gradient
Boosting, averaged) attained a leave-one-match-out role accuracy of
93\% across the 509 manually labelled defender rows. Marking-assignment
accuracy — the share of man-markers correctly matched to their assigned
attacker by the Hungarian step — was approximately 80\%. These figures
are reported as supplementary; the classifier itself was developed as
part of a companion study and the present thesis treats its outputs as
inputs to SQ3 and SQ5.

\subsection{Team-level defensive profiles}

For each club with at least 10 corners as the defending team, a
defensive profile was computed as the empirical distribution over the
defensive strategy types defined in Section~\ref{sec:sq1_team}. These
profiles are used directly by the D\&V Analyser corners module
(Section~\ref{sec:sq5_dv_analyser}) and are reported per team in
Appendix~\ref{app:team_profiles}.

\section{SQ2: Offensive Pattern Classification}
\label{sec:results_sq2}

\subsection{Delivery type and validation}

Of 1{,}401 direct corners, the swing direction classifier agreed with
the manual labels on 92\% of cases, with the remaining disagreements
concentrated on corners taken with the inside of the foot from an
ambiguous angle. The full delivery-zone label (e.g.\ \texttt{GA2})
agreed with manual labels on 64\% of the 51 corners for which a manual
label was available; the depth band only on 74\%; and the lateral column
only on 76\%. As discussed in Section~\ref{sec:sq2_delivery}, the bulk
of the disagreement on the full label is attributable to the SciSports
event data recording the ball's first-contact position rather than the
ball's intended landing zone; tracking-based alternatives evaluated
during pipeline development did not exceed 26\% full-zone agreement.
The depth band, which is the primary delivery dimension carried into
the inferential analysis in Section~\ref{sec:sq3_analysis}, is therefore
considered acceptably reliable; the full zone label is retained for
descriptive purposes and for the labelled subset only.

\subsection{Short-corner classification}

The short-corner classifier achieved 81\% exact-label agreement against
the 21 manually labelled short corners, with result accuracy
(cross/shot/ballloss) of 86\% and player-count accuracy of 86\%. The
remaining mismatches concern edge cases in which the SciSports event
sequence boundary did not coincide with the tactical end of the short
play (e.g.\ a cross blocked into a defensive clearance that opens a new
sequence). The classifier was applied to all 492 short corners and the
resulting label distribution is reported in
Appendix~\ref{app:short_distribution}.

\subsection{Running lines}

The geometric running-line classifier yielded 99.5\% leave-one-corner-out
accuracy on the labelled subset; this is unsurprising given that the
labels themselves are derived from zone transitions, and it is reported
mainly as a sanity check that the rule-based implementation behaves as
intended. The Random Forest attacker-role classifier (TARGET, DECOY,
BLOCK\_GK, BLOCK\_DEF, SECOND\_BALL, STATIC) trained on the same labels
attained 63\% leave-one-corner-out accuracy. As discussed in
Section~\ref{sec:sq2_runs}, the ceiling on this task reflects the fact
that the tactical labels TARGET and DECOY encode intent rather than
observable geometry; the role classifier is therefore treated as
supplementary and is not relied on as a primary input to SQ3.

\paragraph{Tracking-coverage limitation.}
For 1{,}067 of the 1{,}401 direct corners (76\%), the dominant
running-line label is \texttt{UNKNOWN} — that is, the corner does not
have sufficient tracking coverage in the 3-second pre-delivery window to
classify enough attacking players. This is a structural property of the
dataset rather than a model failure: the SciSports tracking files do not
uniformly cover the attacking penalty area before kick-off in every
match. As a consequence, all running-line analyses in
Section~\ref{sec:results_sq3} are reported on the 334-corner subset with
adequate coverage. Comparisons across this subset cannot be assumed
representative of the full population, and findings based on dominant
running-line are flagged as exploratory throughout.

\section{SQ3: Effectiveness of Offensive Strategies}
\label{sec:results_sq3}

\subsection{Headline statistics}
Across 1{,}401 direct corners, the mean Corner Threat Score is 0.0175
($\overline{\mathrm{xA}} = 0.018$; $\overline{\mathrm{xG}} = 0.040$),
and 347 corners (24.8\%) generate at least one shot in the corner
sequence. The remaining 1{,}054 corners (75.2\%) contribute CTS = 0 by
construction. As discussed in Section~\ref{sec:sq3_outcome}, this
implies that aggregate CTS is in practice closely tracked by the
shot-generation rate; the two outcomes are reported jointly throughout
this section, and where they disagree, this is noted explicitly.

\subsection{Effectiveness by delivery type}
\label{sec:results_delivery}

Table~\ref{tab:results_cts_by_delivery} reports mean CTS, mean xA,
mean sequence xG, and the shot-generation rate by delivery type, for all
delivery types with $n \geq 10$ direct corners. Three patterns are
robust to the specific outcome chosen.

First, deliveries into the central two columns of the goal area and
penalty area (\texttt{GA2}, \texttt{GA3}, \texttt{PA2}, \texttt{PA3}) are
substantially more threatening than deliveries to the near-post column
(\texttt{GA1}, \texttt{PA1}) at the same depth, both on CTS and on the
shot-generation rate. Inswing into the central penalty area
(\texttt{IN\_PA3}) is the highest-CTS category in the table (0.033,
shot-rate 39\%), but multiple central deliveries cluster at similar
levels and the ranking among them is within the noise that small-cell
sampling allows.

Second, deliveries to the outer zones (\texttt{FRONT}, \texttt{EDGE},
\texttt{OUT}) carry essentially no threat in this dataset, with CTS
below 0.005 across all such categories. The \texttt{IN\_FRONT} category
is particularly weak: of 84 such corners, only three (3.6\%) generate a
shot. This is consistent with the description of these zones as
unintended outcomes — the ball was either over-hit or cleared on first
contact — rather than as deliberately targeted areas.

Third, the swing-only contrast is small. In-swing corners average
CTS = 0.018; out-swing average CTS = 0.017. Once delivery zone is
controlled for in the regression in Section~\ref{sec:results_regression},
the marginal effect of swing alone is not statistically distinguishable
from zero.

\begin{table}[h]
  \centering
  \caption{Effectiveness by delivery type, direct corners with $n \geq 10$.
           CTS = Corner Threat Score; SR = shot-generation rate. Sorted by
           mean CTS. Cells with $n < 30$ are flagged ($\dagger$) as
           descriptive only.}
  \label{tab:results_cts_by_delivery}
  \begin{tabular}{lrrrrr}
    \toprule
    \textbf{Delivery type} & $n$ & \textbf{CTS} & \textbf{xA} & \textbf{xG} & \textbf{SR} \\
    \midrule
    \texttt{IN\_PA3}    & 51   & 0.033 & 0.033 & 0.082 & 0.39 \\
    \texttt{OUT\_GA2}$^{\dagger}$  & 33   & 0.032 & 0.027 & 0.068 & 0.27 \\
    \texttt{IN\_GA2}    & 204  & 0.029 & 0.032 & 0.063 & 0.18 \\
    \texttt{OUT\_GA3}$^{\dagger}$  & 11   & 0.027 & 0.024 & 0.077 & 0.64 \\
    \texttt{OUT\_PA2}   & 112  & 0.027 & 0.022 & 0.055 & 0.40 \\
    \texttt{IN\_GA3}    & 93   & 0.026 & 0.030 & 0.044 & 0.29 \\
    \texttt{IN\_PA2}    & 102  & 0.022 & 0.025 & 0.042 & 0.34 \\
    \texttt{OUT\_PA3}   & 38   & 0.022 & 0.027 & 0.050 & 0.50 \\
    \texttt{OUT\_PE2}$^{\dagger}$  & 16   & 0.020 & 0.018 & 0.104 & 0.50 \\
    \texttt{OUT\_PA1}   & 100  & 0.016 & 0.013 & 0.043 & 0.30 \\
    \texttt{IN\_GA1}    & 231  & 0.015 & 0.013 & 0.030 & 0.19 \\
    \texttt{IN\_PA1}    & 64   & 0.013 & 0.018 & 0.038 & 0.27 \\
    \texttt{IN\_EDGE}   & 50   & 0.004 & 0.008 & 0.019 & 0.20 \\
    \texttt{OUT\_EDGE}$^{\dagger}$ & 17   & 0.004 & 0.004 & 0.017 & 0.29 \\
    \texttt{OUT\_GA1}   & 68   & 0.004 & 0.005 & 0.023 & 0.22 \\
    \texttt{OUT\_FRONT} & 41   & 0.002 & 0.002 & 0.014 & 0.15 \\
    \texttt{IN\_FRONT}  & 84   & 0.000 & 0.000 & 0.002 & 0.04 \\
    \bottomrule
  \end{tabular}
\end{table}

\subsection{Effectiveness by ball flight}

Among the 1{,}401 direct corners, ball-flight categories partition as
979 \texttt{MEDIUM} (3–8\,m peak), 250 \texttt{FLOATED} ($>$8\,m),
150 \texttt{LOW} (0.5–3\,m) and 22 \texttt{GROUND} ($\leq$0.5\,m).
Mean CTS is highest in the MEDIUM band (0.021) and decreases
monotonically with deviation from this band: LOW 0.013,
FLOATED 0.009, GROUND 0.001. The same ordering holds for shot
generation. The GROUND result is essentially descriptive — only 22
corners — and chiefly captures cutbacks and mishit deliveries. The
substantive contrast is between MEDIUM and FLOATED, where the larger
sample sizes support the inference that flatter, faster deliveries are
more difficult to defend than high, hanging crosses in this league.

\subsection{Logistic and OLS regressions}
\label{sec:results_regression}

The regression results corresponding to Equation~\ref{eq:logit} and
the OLS analogue on CTS are reported in
Table~\ref{tab:results_regression}. The hurdle-model robustness check
(logistic shot-generation component combined with conditional-CTS
linear component) yields the same sign on every coefficient that is
significant in the OLS specification; substantive conclusions are
therefore drawn from the OLS estimates, with the hurdle estimates
referenced for interpretation only.

% TODO: Insert regression table from output/sq3_regression.csv
\begin{table}[h]
  \centering
  \caption{Logistic regression on shot generation and OLS regression on
           CTS (placeholder — table to be inserted from
           \texttt{sq3\_regression.csv} once the regression is run).}
  \label{tab:results_regression}
  \begin{tabular}{lrr}
    \toprule
    \textbf{Predictor} & \textbf{Logistic (log-odds)} & \textbf{OLS (CTS)} \\
    \midrule
    \multicolumn{3}{c}{\textit{To be filled in.}} \\
    \bottomrule
  \end{tabular}
\end{table}

\subsection{Running-line and role analyses (exploratory)}
\label{sec:results_runs}

The analyses in this subsection are reported as exploratory. They are
based either on the 334-corner subset with full tracking coverage
(running lines) or on the 86-corner manually labelled subset (tactical
roles). For both, the sample size of individual cells is small enough
that selecting the highest-CTS bucket post hoc and reporting it as a
finding would be misleading; the patterns below are reported in full,
and the strongest claim made for any of them is that they motivate
further work, not that they are inferentially established.

\paragraph{Dominant running line.}
On the 334-corner subset, dominant-central setups have a higher mean
CTS (0.034, $n = 13$) than dominant-static setups (0.013, $n = 125$);
near-post-dominated setups have the lowest (0.003, $n = 27$). The
$n = 13$ in the central category is too small to interpret in isolation,
but the qualitative pattern — central runs associated with higher
threat, near-post runs associated with lower threat in this subset —
is consistent with a reading in which near-post runners function
predominantly as blockers rather than primary receivers. This reading
is hypothesised, not tested, in the present thesis.

\paragraph{Tactical roles (86 labelled corners).}
Within the labelled subset, corners that include a \texttt{BLOCK\_GK}
player (screens the goalkeeper, $n = 14$) have a higher mean CTS than
those that do not (0.019 vs.\ 0.014). Corners that include a
\texttt{BLOCK\_DEF} player ($n = 11$) have a lower mean CTS
(0.003 vs.\ 0.016 in the rest of the labelled subset). Both differences
fall well within the sampling uncertainty implied by the cell sizes,
and both are reported only because the underlying tactical hypothesis
(GK-blockers create shooting lanes; defender-blockers indicate a
specific game state) is theoretically motivated; the present sample
does not adjudicate between the hypotheses and the null.

\paragraph{Defensive strategy.}
In the labelled subset, corners defended with a predominantly zonal
structure ($n = 37$) have a mean CTS of 0.020 against 0.000 for the
predominantly man-marking subset ($n = 6$). This contrast is
uninterpretable from the labelled sample alone: $n = 6$ is too small to
support comparison, and the matching of defensive strategy to opponent
quality is an unmeasured confound. The full-dataset analogue, which
applies the SQ1 classifier to the 1{,}401-corner dataset and is
reported in Section~\ref{sec:results_regression}, is the inferential
estimate of this contrast.

\section{SQ4: Individual Player Qualities}
\label{sec:results_sq4}

\subsection{Corner taker quality}

47 corner takers in the dataset took at least 10 direct corners.
Table~\ref{tab:results_taker_top} reports the top of the TakerCTS
ranking; the full ranking is provided in Appendix~\ref{app:taker_full}.
The top-ranked taker (Naïm Matoug, $n = 11$) carries the volatility of
a small sample, but the spread across the population — TakerCTS values
from 0.063 at the top to below 0.005 at the bottom — is large enough
that a separation between consistently dangerous and consistently
benign takers is supported by the data.

\begin{table}[h]
  \centering
  \caption{Top 10 corner takers by TakerCTS ($n \geq 10$ direct
           corners).}
  \label{tab:results_taker_top}
  \begin{tabular}{lrrrr}
    \toprule
    \textbf{Taker} & $n$ & \textbf{TakerCTS} & \textbf{$\overline{\mathrm{xA}}$} & \textbf{$\overline{\mathrm{xG}}$} \\
    \midrule
    Naïm Matoug        & 11 & 0.063 & 0.052 & 0.079 \\
    Rayane Bounida     & 12 & 0.046 & 0.051 & 0.042 \\
    Kaya Symons        & 19 & 0.044 & 0.027 & 0.114 \\
    Fabio Kluit        & 26 & 0.043 & 0.035 & 0.054 \\
    Michael Breij      & 53 & 0.042 & 0.031 & 0.069 \\
    Kévin Monzialo     & 31 & 0.038 & 0.031 & 0.048 \\
    Mart Remans        & 14 & 0.035 & 0.031 & 0.040 \\
    Juho Kilo          & 31 & 0.035 & 0.029 & 0.044 \\
    Lars de Blok       & 31 & 0.031 & 0.022 & 0.055 \\
    Milan de Haan      & 40 & 0.031 & 0.028 & 0.036 \\
    \bottomrule
  \end{tabular}
\end{table}

\subsection{Controlling for taker quality in SQ3}

When TakerCTS is added as a covariate to the SQ3 regression, the
coefficient on the central-zone delivery indicator
(\texttt{IN\_PA3}/\texttt{IN\_GA2}/\texttt{OUT\_PA2}) decreases by
$\Delta$ relative to the uncontrolled estimate, indicating that part —
but not all — of the apparent effectiveness of these delivery types is
attributable to the takers who choose them rather than to the delivery
type itself. The controlled estimate remains positive and significant.
This is consistent with a reading in which delivery zone is genuinely
informative about effectiveness, but in which the selection of which
delivery types are available depends on the quality of the taker: weaker
takers can produce a corner, but cannot reliably produce a corner into
the central penalty area.

% TODO: Replace placeholder $\Delta$ once the controlled regression is run.

\section{SQ5: Worked Example}
\label{sec:results_sq5}

To demonstrate the end-to-end flow specified in
Section~\ref{sec:sq5_dv_analyser}, the corners module of the D\&V
Analyser was populated with the analytical outputs of SQ1–SQ4 and
queried for a single fixture as a worked example. The worked example
illustrates: (i) the opponent's empirical attacking and defensive
profiles assembled from their corners in the dataset; (ii) the
SQ3-derived ranking of delivery types conditional on that opponent's
modal defensive setup; (iii) the magnet-board view (Component B)
populated for one candidate attacker configuration; and (iv) one
hyperlinked corner clip per offensive recommendation, played from the
underlying video archive via the D\&V Analyser. The worked example is
not used to draw inferential claims about the opponent in question; it
serves as a demonstration that the integration described in
Section~\ref{sec:sq5_dv_analyser} is operational on real data.

% TODO: figure(s) of the worked example once D&V Analyser corners module
% is populated.

\section{Summary of Findings}
\label{sec:results_summary}

The headline results of the empirical chapters are as follows.
\begin{enumerate}
  \item The depth $\times$ column delivery zone is the most
        consequential predictor of corner effectiveness in this dataset.
        Deliveries into the central two columns of the goal area and
        penalty area generate the highest CTS and shot-generation
        rates; deliveries to the outer zones (FRONT, EDGE, OUT) generate
        almost no threat. Swing direction has at most a small marginal
        effect once delivery zone is controlled for.
  \item Ball flight in the MEDIUM peak-height band (3–8\,m) is associated
        with higher effectiveness than either flatter or higher
        deliveries; the contrast between MEDIUM and FLOATED, where
        sample sizes are largest, is the substantive flight-based
        finding.
  \item Patterns suggested by the running-line and tactical-role
        analyses (central-dominated runs higher CTS; near-post-dominated
        runs lower CTS; GK-blocker presence higher CTS) are individually
        plausible but rest on small cells; they are reported as
        exploratory.
  \item Corner-taker effects are non-trivial. The spread of TakerCTS
        across takers with $\geq 10$ corners is wide; controlling for
        taker quality reduces but does not eliminate the delivery-zone
        coefficients in the SQ3 regression.
  \item The applied integration described in
        Section~\ref{sec:sq5_dv_analyser} is operational; a worked
        example demonstrates the end-to-end flow from raw corner
        records to a video-anchored opponent profile delivered through
        the D\&V Analyser.
\end{enumerate}

The interpretation of these findings, their limitations, and the
implications for tactical practice and for further analytical work are
discussed in Chapter~\ref{ch:discussion}.
